\documentclass[11pt]{article}
\usepackage{amsmath,amssymb,amsthm,hyperref,geometry}\geometry{margin=1in}
\newtheorem{theorem}{Theorem}\newtheorem{proposition}{Proposition}
\title{Boundary-admitted joint occurrence unfolding}
\date{15 November 2026}\begin{document}\maketitle
\paragraph{Status.} Scoped theorem with written proofs. This joins local effect admission, Cartesian occurrence content and the existing retained guarded unfolding. The joint next-state body is constructed, not assumed. It is a scoped general theorem for monadic effect contexts and polynomial occurrences; probability and nondeterminism are separate realizations. Shared boundary laws are substantive hypotheses, not a claim that arbitrary hidden interference disappears.

\section{Component systems and retained overlap}
Let a finite nonempty tree index state sets $X_v$ and overlap spans $X_v\xrightarrow{b_{ve}}B_e\xleftarrow{b_{we}}X_w$. Let $X$ be their full compatible-tuple set with projections $h_v:X\to X_v$. Let $S$ be a finitary Cartesian polynomial monad. Initially let $T$ merely be a Set functor weakly preserving pullbacks.

At each component state $x_v$ supply a raw witness set $W^v_{x_v}$, effect $p^v_{x_v}\in T W^v_{x_v}$ and occurrence body
\[
 \phi_v(x_v,-):W^v_{x_v}\to S X_v.
\]
At a fixed compatible joint state $x\in X$, impose the edge equations
\[
 T(Sb_{ve}\circ\phi_v(h_vx,-))(p^v_{h_vx})
 =T(Sb_{we}\circ\phi_w(h_wx,-))(p^w_{h_wx})
 \quad\text{in }T(SB_e).
\tag{E}
\]
These equations compare the full boundary occurrence effects, including addressed shape, not only an average terminal observation.

A joint step at $x$ means a raw set $Z_x$, $q_x\in TZ_x$, maps $u_{v,x}:Z_x\to W^v_{h_vx}$ and a body $\phi(x,-):Z_x\to SX$ with
\[
 Tu_{v,x}(q_x)=p^v_{h_vx},\qquad
 Sh_v\,\phi(x,z)=\phi_v(h_vx,u_{v,x}z).
\tag{N}
\]
Thus each component projection is an actual native context/body square.

\begin{theorem}[Exact local admission and constructed body]
A joint step exists if and only if (E) holds. Under (E), it can be chosen with raw carrier the full set of component witness tuples whose boundary occurrences match on every edge. The effect and the joint occurrence body on that carrier then satisfy (N).
\end{theorem}
\begin{proof}
Necessity follows from $b_{ve}h_v=b_{we}h_w$ and the two native equations. For sufficiency apply weak-pullback tree admission to the raw witness fibres and edge maps $Sb_{ve}\phi_v$, using (E). This supplies an effect $q_x$ on their compatible-tuple set $Z_x$ with the prescribed component witness marginals.

The component bodies define a map from $Z_x$ to the compatible-tuple set of the diagram of sets $SX_v$ over $SB_e$. Because $S$ preserves pullbacks and the tree is finite and nonempty, repeated leaf decomposition identifies that set canonically with $SX$. This identification constructs $\phi(x,-)$ on every raw matching tuple, including tuples outside any effect support. Its component equations are exactly the second part of (N). No conditional disintegration or canonical selection of the effect lift has been used.
\end{proof}

Suppose (E) holds for every $x\in X$. Choosing a lift $q_x$ at each state, with ordinary Set choice or an explicitly supplied family, gives a joint system. Choices remain retained data. Neither local equality nor the theorem identifies different couplings with the same marginals. Conversely any joint system with those native projections forces the equations at each state.

The equations also express a precise information boundary: although the selected joint effect can depend on the whole $x$, its component marginal must equal the specified component law at $h_vx$. If an independently proposed joint law changes that prescribed component law through unrecorded state, it fails the native marginal equation (N). Boundary equality (E) alone need not detect this: a constant boundary can hide such changes. The existence theorem constructs some compatible joint system; it does not certify an unspecified physical candidate as that system. One must enlarge the retained interface or alter the component specification; averaging over selected joint histories is not a substitute.

\section{Guarded recursive assembly}
Now suppose $T$ is a commutative affine monad and its underlying functor has the weak-pullback property. Use the existing addressed unfolding convention: at each finite frontier choose the finite product of its current-state effects, retain every raw witness and occurrence address, and advance through the supplied $S$ body. Different frontier occurrences use that specified Fubini construction. Arbitrary cross-occurrence dependence is not silently added by this convention; relevant memory may be retained in the state or require a different context construction.

\begin{theorem}[Joint native prefix system]
The chosen joint system has compatible finite-depth raw/effect prefixes. Each $h_v,u_v$ extends to a native map of its entire retained prefix system to the $v$th component system. These maps preserve the complete addressed occurrence bodies and push the joint prefix effect to the component prefix effect at every depth. They commute with truncation. Thus the joint coupling persists through arbitrary finite guarded unfolding without checking a new global marginal equation.
\end{theorem}
\begin{proof}
At depth zero use the initial state projections. For a fixed addressed frontier, the second equation in (N) identifies component and joint occurrence shapes and their positions; equal labels do not identify addresses. The map on newly chosen frontier witnesses is the product of the actual $u_v$ maps at the corresponding current joint states. Naturality of the commutative monad's finite Fubini products and the first equation in (N) push the joint frontier effect to the specified component frontier effect. Binding this conditional extension onto the retained preceding history gives the component prefix effect by monad naturality. The body equation advances corresponding child states according to $h_v$, so the induction is well-typed at the next depth. All raw histories are retained, not only final flattened labels. When the displayed body is flattened to $SX$ by the monad multiplication, naturality of that multiplication gives its projected body equation as well; the raw hierarchy is not identified with that flattened endpoint.

Truncation forgets only the last layer. Affinity gives the effect equation for this forgetting map, as in the prior guarded native unfolding theorem, and the raw projection maps plainly commute with it. This completes the induction on depth. The constructed joint occurrence body, rather than a separately assumed global correspondence, is what supplies the induction's state compatibility.
\end{proof}

Compatible raw limits may be taken separately. No infinite $T$-effect or measure is inferred. If every $u_{v,x}$ is surjective on raw witnesses at every reachable joint state, then every component raw unfolding from $h_vx$ lifts from that specified joint state $x$, by successive choices at each addressed node. Finitary branching gives countably many finite-depth nodes; ordinary dependent choice suffices. This is an extra local-surjectivity hypothesis, not a consequence of (E): zero-weight or disallowed witnesses can lack compatible partners. No fairness or scheduler implementability is implied.

\section{Integration and limits}
The construction gives a star diagram of actual native occurrence systems. The existing diagram-native interaction transports its native structure; the present theorem additionally derives the root system from boundary equations. A finite tree of component overlaps is not replaced by an arbitrary cyclic wiring theorem. Pairwise consistency around a cycle can fail global admission.

Finite-support distributions and nonempty powerset are different instances of commutative affine weak-pullback monads. In the former, joint effects retain statistical dependence; in the latter they retain allowed joint relations. The theorem is stated at the monadic/native level, not defined by either model. It relies on the established weak-pullback tree-admission and polynomial guarded-unfolding machinery, and makes no literature-priority claim. Its added task is their explicit integration: local boundary equations construct a joint body whose native squares justify every guarded prefix comparison.

This is a guarded recursive joint-assembly theorem, not a general trace, a universal feedback replacement, or an arbitrary shared-interference interactad. Boundary compatibility must hold at all admitted joint states, and the frontier effect policy is the stated one.
\end{document}
